\section{$k$-forme}

Una funzione $\lambda\colon (\RR^n)^k$, 
$(v_1,\dots, v_k)\mapsto \lambda[v_1,\dots, v_k]$
con $v_1,\dots, v_k \in \RR^n$
si dice essere \emph{multilineare} se ogni 
$v_k \mapsto \lambda[v_1,\dots, v_k]$ 
è lineare e se scambiando due ingressi $v_k$, $v_j$ 
il risultato cambia segno.

\begin{definition}[applicazioni multilineari alternanti]
Sia $V$ è uno spazio vettoriale reale.
Denotiamo con $\Lambda^k(V)$
l'insieme delle applicazioni multilineari alternanti.

Cioè $\lambda \in \Lambda^k(V)$ se 
$\lambda\colon (V)^k\to \RR$ e
\begin{enumerate}
\item è lineare nella prima entrata:
\[
    \lambda[tv + s v', v_2, \dots, v_k] 
    = t\lambda[v,v_2,\dots,v_k]
    + s \lambda[v',v_2,\vdots, v_k],
\]
\item è alternante:
\[
    \lambda[v_{\sigma(1)}, \dots v_{\sigma(k)}] = 
    (-1)^\sigma \lambda[v_1,\dots, v_k]
    \]
per ogni $\sigma$ permutazione degli indici $\ENCLOSE{1,\dots,k}$
dove $(-1)^\sigma$ è il segno della permutazione.
\end{enumerate}
\end{definition}

La definizione assiomatica del determinante 
ci dà questo importante esempio.
\begin{example}
Sia $\lambda\in \Lambda^n(\RR^n)$. 
Allora esiste $c\in \RR$ tale che 
\[ 
    \lambda[v_1, \dots, v_n] = c \det \begin{pmatrix}
    v_1, \dots, v_n
    \end{pmatrix}.
\]
\end{example}

\begin{definition}[$k$-forme differenziali]
Sia $A$ un aperto di $\RR^n$.
Una $k$-forma differenziale $\omega$ è una funzione 
$\omega\colon A \to \Lambda^k(\RR^n)$.
L'insieme delle $k$-forme su $A$ si denota con $\Omega^k(A)$.
\end{definition}

Abbiamo già considerato lo spazio delle $1$-forme e abbiamo 
visto che possono essere rappresentate nella forma 
\[
  \omega = \xi_1 dx_1 + \dots + \xi_n dx_n, \qquad \omega(x)[v] = \xi(x)\cdot v
\]
con $\xi\colon A \to \RR^n$.

Le $0$-forme sono banalmente le funzioni scalari, in quanto $\Lambda^0(\RR^n)= \RR$
e dunque $f \in \Omega^k(A)$ significa $f\colon A \to \RR$.

Visto che $\Lambda^n(\RR^n)$ contiene solamente i multipli del determinante, 
anche le $n$-forme in $\RR^n$ si rappresentano con una funzione $f\colon A\to\RR$, 
\[
 \omega(x)[v_1,\dots,v_n] = f(x) \det(v_1, \dots, v_n),\qquad f(x) = \omega(x)[e_1,\dots,e_n].
\]

I casi intermedi sono più complicati da capire.
Se $\lambda\in \Lambda^k(\RR^n)$ sappiamo che se scambiamo due vettori 
in $\lambda[v_1,\dots,v_n]$ il risultato cambia segno. 
In particolare se ci sono due vettori uguali $v_k=v_j$ con $k\neq j$, il risultato 
è nullo, in quanto scambiandoli dovrebbe cambiare segno ma si ottiene lo stesso risultato.
Dunque $\lambda$ è non nullo solamente sulle $k$-uple di vettori $(v_1,\dots,v_k)$ 
che non hanno ripetizioni. E il valore su una $k$-upla determina il valore su tutte le sue 
permutazioni. 
Si capisce dunque che la $\dim \Lambda^k(\RR^n) = {k \choose n}$.

Infatti abbiamo già visto che $\dim \Lambda^0(\RR^n)=\dim \Lambda^n(\RR^n) = 1$.

Proviamo a capire come è fatta $\Omega^2(\RR^n)$. 
Consideriamo $dx_1,\dots, dx_n \in \Omega^1(\RR^n)$. 
Sono definiti in modo che $dx_k(p)[v] = v_k$.
Definiamo allora 
$dx_k \wedge dx_j\in \Omega^2(\RR^n)$ tramite
\[
    (dx_k\wedge dx_j)(p) [v,w] = dx_k(p)[v] dx_j(p)[w] - dx_k(p)[w] dx_j(p)[v]
      = v_k w_j - v_j w_k. 
\]
E' chiaro che questa è una forma bilineare alternante.
E, dalla discussione precedente, è chiaro che ogni forma bilineare alternante 
è combinazione lineare di queste.
Infatti se $\omega \in \Omega^2(\RR^n)$ 
posto $\xi_{jk}(x) = \omega(x)[v_j,v_k]$ per ogni $1\le j<k\le n$
si ha 
\[
  \omega(x) = \sum_{1\le j < k \le n} \xi_{jk}(x) dx_j\wedge dx_k.
\] 

\begin{example}[$2$-forme in $\RR^3$]
Sia $\omega \in \Omega^2(\RR^3)$. 
Denotiamo con $x,y,z$ le tre coordinate di $\RR^3$ (invece che $x_1,x_2,x_3$).
Allora
\[
  \omega = \xi_{12} dx\wedge dy + \xi_{13} dx\wedge dz + \xi_{23} dy \wedge dz.
\] 
Notiamo che si ha $dy\wedge dx = - dx\wedge dy$ e lo stesso per le altre coordinate.
Dunque si potrebbe anche scrivere 
\[
 \omega = \xi_{23} dy \wedge dz - \xi_{13} dz\wedge dx + \xi_{12} dx\wedge dy.
\]
\end{example}

\section{differenziale di una forma}

Possiamo definire il \emph{differenziale} $d\omega$ di una forma differenziale $\omega$, 
definendo il differenziale delle forme \emph{elementari}
\[
  d( f dx_{j_1} \wedge \dots  \wedge dx_{j_k})
  = \sum_i \frac{\partial f}{\partial x_i} x_i \wedge dx_{j_1} \wedge \dots \wedge dx_{j_k}.
\]
Lo si capisce con un esempio.

\begin{example}[calcolo del differenziale]
Sia $\omega\in \Omega(\RR^3)$, $\omega = f dx \wedge dy - g dy\wedge dz$ con $f,g\colon \RR^3\to \RR$.
Allora 
\begin{align*}
  d\omega 
  &= df \wedge dx \wedge dy - dg \wedge dy \wedge dz\\
  &= (f_x dx + f_y dy + f_z dz) \wedge dx \wedge dy - (g_x dx + g_y dy + g_z dz)\wedge dy \wedge dz\\
  &= f_z dz\wedge dx \wedge dy - g_x dx \wedge dy \wedge dz \\
  &= (f_z - g_x) dx \wedge dy \wedge dz.
\end{align*}
Ricordiamo che termini come $dx\wedge dx \wedge dy$ o $dy \wedge dx \wedge dy$ sono nulli in quanto 
si ripete due volte lo stesso differenziale. 
Invece $dz\wedge dx \wedge dy = dx \wedge dy \wedge dz$ perché si tratta di una permutazione 
pari (due trasposizioni). Si avrebbe invece $dx\wedge dz\wedge dy = - dx\wedge dy\wedge dz$.
\end{example} 

Abbiamo già visto sulle $1$-forme che se $\omega = df$ (cioè $\omega$ è esatta) 
allora $d\omega = d(df) = 0$ (cioè $\omega$ è chiusa).

In generale si può dimostrare che vale:
\[
  d(d\omega) = 0.
\]

\section{pull-back}

Se $\phi\colon A\subset \RR^m \to \RR^n$ è una funzione differenziabile e $\omega\in \Omega^k(\phi(A))$
possiamo definire il pull-back $\phi^* \omega \in \Omega^k(A)$ tramite la formula
\[
  (\phi^* \omega)(x)[v_1,\dots,v_k] = \omega(\phi(x))[d\phi(x)[v_1],\dots, d\phi(x)[v_k]].
\]

Osserviamo che per definizione l'integrale di una $k$-forma differenziale 
su una $k$-superficie parametrizzata da $\phi\colon [0,1]^k \to \RR^n$
si definisce come l'integrale del pull-back della $k$-forma:
\[ 
  \int_\phi \omega \defeq \int_{[0,1]^k} \phi^* \omega
  = \int_{[0,1]^k} f(\vec x) d\vec x
\]
dove $f\colon [0,1]^k \to \RR$ è la funzione che rappresenta $\phi^* \omega$
in quanto $\phi^* \omega$ è una $k$-forma su $[0,1]^k$ e dunque si rappresenta 
con una funzione scalare $f$: $\phi^* \omega(\vec x) = f(\vec x) dx_1 \wedge \dots \wedge dx_k$.

In particolare, visto che il pull-back di una composizione è ovviamente 
la composizione dei pull-back, se $\vec u \colon [0,1]^k\to \RR^m$, 
$\phi\colon \RR^m\to \RR^n$ e $\omega\in \Omega^k(\RR^n)$
si ha 
\[
   \int_{\vec u} \phi^* \omega = \int_{\phi \circ \vec u} \omega 
\]
in quanto entrambi si riconducono a
\[
  \int_{[0,1]^k} (\phi \circ \vec u)^* \omega.
\]

In generale si può dimostrare che 
\[
    \phi^* (d\omega) = d(\phi^* \omega).
\]

\section{$k$-forme in $\RR^2$}
Se $\omega\in \Omega^1(A)$, con $A$ aperto di $\RR^2$, si ha 
\[
  \omega(x,y) = a(x,y) dx + b(x,y) dy
\]
con $a,b\colon A \to \RR$.

Calcoliamo il differenziale di $\omega$. 
Si ha 
\begin{equation}
  \label{eq:differenziale_bidimensionale}
  \begin{aligned}
 d \omega 
  &= da\wedge dx + db\wedge dy\\
  &= \enclose{\frac{\partial a}{\partial x} dx + \frac{\partial a}{\partial y}dy}\wedge dx 
   + \enclose{\frac{\partial b}{\partial x} dx + \frac{\partial b}{\partial y}dy}\wedge dy \\
  &= \frac{\partial a}{\partial y} dy\wedge dx 
   + \frac{\partial b}{\partial x} dx \wedge dy \\
  &= \enclose{\frac{\partial b}{\partial x} 
    - \frac{\partial a}{\partial y}} dx \wedge dy.
\end{aligned}
\end{equation}


\section{$k$-forme in $\RR^3$}

\begin{definition}[prodotto vettore]
Siano $v,w\in \RR^3$.
Definiamo il \emph{prodotto vettore} $v\times w\in \RR^3$ 
come l'unico vettore tale che per ogni $u\in \RR^3$ si abbia
\[
  u \cdot (v\times w) = \det\enclose{u, v, w} \qquad\text{(prodotto triplo)}.
\]
Questa uguaglianza definisce univocamente $v\times w$ in quanto 
il lato destro è una applicazione lineare in $u$ e dunque 
c'è un unico vettore che la rappresenta mediante prodotto scalare.
\end{definition}

In coordinate, se $\xi = v\times w$ con $v=(v_1,v_2,v_3)$, $w=(w_1,w_2,w_3)$
si ha dunque $\xi=(\xi_1,\xi_2,\xi_3)$ 
con
\begin{align*}
  \xi_1 &= e_1 \cdot \xi = \det\enclose{e_1,v,w} = v_2 w_3 - v_3 w_2,\\ 
  \xi_2 &= e_2 \cdot \xi = \det\enclose{e_2,v,w} = v_3 w_1 - v_1 w_3,\\
  \xi_3 &= e_2 \cdot \xi = \det\enclose{e_3,v,w} = v_1 w_2 - v_2 w_1.
\end{align*}

Più espressivamente si può scrivere
\[
  \xi = \det\begin{pmatrix}
    e_1 & v_1 & w_1 \\ e_2 & v_2 & w_2 \\ e_3 & v_3 & w_3
  \end{pmatrix}
  = (v_2 w_3 - v_3 w_2) e_1 + (v_3 w_1 - v_1 w_3) e_2 + (v_1 w_2 - v_w w_1) e_3. 
\]

\begin{definition}[rotore]
Sia $A$ un aperto di $\RR^3$ e $\xi \colon A \to \RR^3$ una funzione differenziabile.
Sia $e_1,e_2,e_3$ la base canonica di $\RR^3$.
Definiamo allora il \emph{rotore} di $\xi$ come il vettore $\rot \xi \in \RR^3$ 
definito da 
\[
   \rot \xi = \nabla \times \xi
    = \det \begin{pmatrix}
    e_1 & \partial_x & \xi_1 \\
    e_2 & \partial_y & \xi_2 \\
    e_3 & \partial_z & \xi_3 
    \end{pmatrix}
    = \begin{pmatrix}
    \frac{\partial \xi_3}{\partial y} - \frac{\partial \xi_2}{\partial z} \smallskip \\
    \frac{\partial \xi_1}{\partial_z} - \frac{\partial \xi_3}{\partial x} \smallskip \\
    \frac{\partial \xi_2}{\partial x} - \frac{\partial \xi_1}{\partial y}
    \end{pmatrix}.
\]

Ricordiamo anche gli operatori \emph{gradiente} e \emph{divergenza}
\[
  \grad f = \nabla f = \begin{pmatrix}
     \frac{\partial f}{\partial x} \smallskip \\
     \frac{\partial f}{\partial y} \smallskip \\
     \frac{\partial f}{\partial z}
  \end{pmatrix}, \qquad 
  \div f = \nabla \cdot \xi = \frac{\partial \xi_1}{\partial x} + \frac{\partial \xi_2}{\partial y} + \frac{\partial \xi_3}{\partial z}.
\]
In tutte queste notazioni abbiamo usato l'operatore \emph{nabla}: $\nabla = (\partial_x, \partial_y, \partial z)$
dove $\partial_x = \frac{\partial}{\partial x}$, $\partial y = \frac{\partial}{\partial y}$, $\partial_z = \frac{\partial}{\partial z}$:
\end{definition}

Abbiamo già visto che una $1$-forma $\omega$ in un aperto $A$ di $\RR^3$ può essere rappresentata 
con un campo vettoriale $\xi\colon A \subset \RR^3 \to \RR^3$ 
in modo tale che si abbia 
\[
  \omega(x,y,z)[v] = \xi(x,y,z)\cdot v.
\]

Se consideriamo $f\colon A\subset \RR^3\to \RR$ una $0$-forma, 
il suo differenziale $\omega = df$ è una $1$-forma, che si rappresenta 
con il vettore $\nabla f$:
\[
df = \frac{\partial f}{\partial x}\, dx 
+ \frac{\partial f}{\partial y} dy\, \frac{\partial f}{\partial z}\, dz.
\]

Se $\omega\in \Omega^1(A)$, possiamo scrivere $\omega$ nella forma 
\[
 \omega = g_1\, dx + g_2\, dy + g_3\, dz
\]
in modo che $\vec g \colon A\to\RR^3$ sia il campo vettoriale che rappresenta $\omega$.
In questo caso si ha 
\begin{align*}
    d\omega &= (\partial_x \vec g_1 dx + \partial_y \vec g_1 dy + \partial z \vec g_3 dz)\wedge dx 
    + (\partial x \vec g_2 dx + \partial y \vec g_2 dy + \partial z \vec g_3 dz)\wedge dy \\
    & \quad + (\partial x \vec g_3 dx + \partial y \vec g_3 dy + \partial z \vec g_3 dz)\wedge dz \\    
    &= (\partial_y \vec g_3 - \partial_z \vec g_2)\, dy \wedge dz
       + (\partial_z \vec g_1 - \partial_x \vec g_3)\, dz\wedge dx 
       + (\partial_x \vec g_2 - \partial_y \vec g_1)\, dx \wedge dy\\
    &= (\rot \vec g)_1 \, dy \wedge dz
       + (\rot \vec g)_2\, dz\wedge dx 
       + (\rot \vec g)_3\, dx \wedge dy.
\end{align*}
Dunque se rappresentiamo $d\omega\in \Omega^2(A)$, tramite 
il campo vettoriale $\xi\colon A\to \RR^3$
definito da 
\[
\xi_1(\vec p) = \omega(\vec p)[e_2,e_3], \qquad
\xi_2(\vec p) = \omega(\vec p)[e_3,e_1], \qquad
\xi_3(\vec p) = \omega(\vec p)[e_1,e_2]
\]
si avrà
\begin{equation}\label{eq:fkjtpd}
  d \omega = \xi_1\, dy \wedge dz + \xi_2\, dz\wedge dx + \xi_3\, dx\wedge dy.
\end{equation}
Abbiamo quindi verificato che se $\xi$ è il campo vettoriale che rappresenta 
la $1$-forma $\omega$, il campo (pseudo)-vettoriale che rappresenta $d\omega$
è $\rot \xi$.

Consideriamo infine $\omega \in \Omega^2(A)$ una $2$-forma 
definita su un aperto $A\subset\RR^3$.
Sappiamo che $\omega$ si rappresenta tramite un campo
vettoriale $\xi \colon A \subset \RR^3 \to \RR^3$  
come in~\eqref{eq:fkjtpd}.
Osserviamo che allora si ha 
\begin{align*}
  \omega[v,w] 
  &= \xi_1 (v_1 w_3 - v_3 w_2) + \xi_2 (v_3 w_1 - v_1 w_3)
   + \xi_3 (v_1 w_2 - v_2 w_1) \\
   &= \det \enclose{\xi, v, w} = \xi \cdot (v\times w).
\end{align*}
Inoltre
\begin{align*}
 d\omega &= (\partial_x \xi_1 dx + \partial_y \xi_1 dy + \partial_z \xi_1 dz) \wedge dy \wedge dz \\
    &\quad + (\partial_x \xi_2 dx + \partial_y \xi_2 dy + \partial_z \xi_2 dz) \wedge dz \wedge dx \\
    &\quad + (\partial_x \xi_3 dx + \partial_y \xi_3 dy + \partial_z \xi_3 dz) \wedge dx \wedge dy \\
    &= (\partial_x \xi_1 + \partial_y \xi_2 + \partial_z \xi_3)\, dx \wedge dy \wedge dz\\
    &= (\div \xi)\, dx \wedge dy \wedge dz.
\end{align*}
Dunque la funzione che rappresenta la $3$-forma $d\omega$ 
è la divergenza del campo vettoriale che rappresenta la $2$-forma $\omega$.

\begin{exercise}
Verifichiamo che vale l'identità
\[
  d (\phi^* \omega) = \phi^* d\omega
\]
quando $\omega$ è una $1$-forma in $\RR^3$.
\end{exercise}

\begin{proof}[Svolgimento.]
Possiamo rappresentare $\omega$ con un campo vettoriale $\xi\colon \RR^3\to \RR^3$:
\[
\omega = \xi_1 dx + \xi_2 dy + \xi_3 dz.
\]
Allora abbiamo già visto che 
\[
  (d\omega)[v,w] = \det \enclose{\rot \xi, v, w}.
\]
Dunque
\[
\phi^*(d\omega)[v,w] = (d\omega)[d\phi[v],d\phi[w]]
 = \det \enclose{\rot \xi, \frac{\partial \phi}{\partial v}, \frac {\partial \phi}{\partial w}}.
\]

D'altro canto
\[
 \phi^* \omega [v] 
 = \omega[d\phi[v]] 
 = \xi \cdot \frac{\partial \phi}{\partial v}.
\]
Ora se $\phi^* \omega$ si rappresenta tramite un 
campo vettoriale $\eta$ in $\RR^2$ (usando le coordinate $t,s$ in $\RR^2$):
\[
  \phi^* \omega = \eta_1 dt + \eta_2 ds
\]
si ha 
\begin{align*}
 \eta_1 &= (\phi^*\omega)[e_1] 
 = \xi \cdot \frac{\partial \phi}{\partial s}\\
 \eta_2 &= (\phi^*\omega)[e_2]
 = \xi \cdot \frac{\partial \phi}{\partial t}
\end{align*}
dunque
\begin{align*}
d(\phi^* \omega) 
 &= \enclose{\frac{\partial \eta_2}{\partial s} - \frac{\partial \eta_1}{\partial t}} ds \wedge dt\\
 &= \enclose{\frac{\partial \xi}{\partial s}\cdot \frac{\partial \phi}{\partial t}
  + \xi \frac{\partial^2 \phi}{\partial s \partial t}}
  - \enclose{\frac{\partial \xi}{\partial t}\cdot \frac{\partial \phi}{\partial s}
  + \xi \cdot \frac{\partial^2 \phi}{\partial t\partial s}
  }\\
 &= \sum_k \enclose{
    \frac{\partial \xi_k}{\partial s} \frac{\partial \phi_k}{\partial t}
    - \frac{\partial \xi_k}{\partial t} \frac{\partial \phi_k}{\partial s}
  }\\
 &= \sum_j \enclose{
    \frac{\partial \xi_k}{\partial x_j} \frac{\partial \phi^j}{\partial s} \frac{\partial \phi_k}{\partial t}
    - \frac{\partial \xi_k}{\partial x_j} \frac{\partial \phi_j}{\partial t}\frac{\partial \phi_k}{\partial s}
  } \\
 &= \sum_{j\neq k} \frac{\partial \xi_k}{\partial x_j} \enclose{
  \frac{\partial \phi_j}{\partial s} \frac{\partial \phi_j}{\partial t}
  - \frac{\partial \phi_j}{\partial t} \frac{\partial \phi_k}{\partial s}}\\
  &= \det\enclose{\rot xi, \frac{\partial \phi}{\partial s}, \frac{\partial \phi}{\partial t}}.
\end{align*}
\end{proof}